\documentclass[11pt,a4paper]{article}
\usepackage[T1]{fontenc}
\usepackage{lmodern}
\usepackage[margin=28mm]{geometry}
\usepackage{amsmath,amssymb,amsthm,booktabs}
\usepackage{microtype}
\usepackage[hidelinks]{hyperref}\hypersetup{pdftitle={A direct-sum law for the Ray--West correction and a statistic for separable permutations},pdfauthor={Xiaojian Zeng},pdfsubject={Permutation patterns and exact combinatorial enumeration}}
\newtheorem{theorem}{Theorem}[section]
\newtheorem{lemma}[theorem]{Lemma}
\newtheorem{corollary}[theorem]{Corollary}
\theoremstyle{remark}
\newtheorem{remark}[theorem]{Remark}
\newcommand{\Ins}{\operatorname{Ins}}
\newcommand{\ind}{\mathbf{1}}
\newcommand{\dec}{\mathrm{dec}}
\newcommand{\inc}{\mathrm{inc}}
\newcommand{\good}{\mathcal{G}}\newcommand{\skewsum}{\mathbin{\mathrm{skew}}}
\title{A direct-sum law for the Ray--West correction\\
and a statistic for separable permutations}
\author{Xiaojian Zeng\thanks{PBC School of Finance, Tsinghua University,
100190, China. Email: \href{mailto:xiaojianzeng@foxmail.com}{xiaojianzeng@foxmail.com}.}}
\date{}

\begin{document}
\maketitle
\begin{abstract}
Ray and West expressed the number of permutations of length $m+2$
containing a given permutation $\pi$ of length $m$ as a polynomial in $m$
minus a correction $j(\pi)\in\{0,\ldots,m-1\}$. We prove that
$j(\alpha\oplus\beta)=j(\alpha)+j(\beta)+T(\alpha)H(\beta)$ for arbitrary
nonempty summands. Here $T$ and $H$ indicate decreasing last and first
sum-indecomposable components, respectively. The proof localizes the
Ray--West insertion witnesses at a direct-sum cut; restriction may require
re-pairing the inserted points. Complementation gives a skew-sum law.
These laws yield an intrinsic decomposition-tree statistic on every
separable permutation and classify equality $j(\pi)=|\pi|-1$ throughout
the sum- and skew-decomposable domains. The problem of expressing $j$
intrinsically for arbitrary permutations remains unresolved by these results.
\end{abstract}

\section{The composition law and its scope}

Write $S_m$ for the permutations of $\{1,\ldots,m\}$ in one-line notation.
A permutation $\tau$ contains $\pi$ if some subsequence of $\tau$ has the
same relative order as $\pi$. Let $g_2(\pi)$ count the \emph{distinct}
permutations in $S_{m+2}$ containing $\pi\in S_m$, rather than their
occurrences. Set
\begin{equation}\label{eq:residual}
 B(m)=\frac{m^4+2m^3+m^2+4m+4}{2},\qquad
 j(\pi)=B(m)-g_2(\pi).
\end{equation}
Ray and West~\cite[Corollary~8.6 and Theorem~8.7]{raywest}
proved $0\le j(\pi)\le m-1$ and described $j$ using insertion tracks.
Vatter~\cite[Problem~3.7]{vatter} asks for a description in terms of
statistics of $\pi$. We answer this question on the separable class and
give a composition law without any separability hypothesis on the summands.
We do not resolve the question for all permutations.

For $\alpha\in S_a$ and $\beta\in S_b$, define
\[
 \alpha\oplus\beta=
 (\alpha(1),\ldots,\alpha(a),a+\beta(1),\ldots,a+\beta(b)),
\]
\[
 \alpha\skewsum\beta=
 (b+\alpha(1),\ldots,b+\alpha(a),\beta(1),\ldots,\beta(b)).
\]
A permutation is \emph{sum-indecomposable} if it is not a direct sum of
two nonempty permutations. Its maximal direct-sum decomposition is its
unique decomposition into sum-indecomposable components. Indeed, the cuts
are exactly the indices $k<m$ with
$\{\pi(1),\ldots,\pi(k)\}=\{1,\ldots,k\}$, and taking all cuts gives
uniqueness. Skew-indecomposability and maximal skew decomposition are
defined analogously. Let
\begin{align*}
 H(\pi)&=\ind\{\text{the first sum component of $\pi$ is decreasing}\},\\
 T(\pi)&=\ind\{\text{the last sum component of $\pi$ is decreasing}\}.
\end{align*}
Singletons are both increasing and decreasing.

\begin{theorem}\label{thm:sum}
For every pair of nonempty permutations $\alpha,\beta$,
\begin{equation}\label{eq:sum}
 j(\alpha\oplus\beta)=j(\alpha)+j(\beta)+T(\alpha)H(\beta).
\end{equation}
More precisely, the Ray--West good gaps internal to either summand
restrict exactly to the good gaps of that summand. The gap at the join
is good if and only if $T(\alpha)H(\beta)=1$.
\end{theorem}

Good gaps are defined in Section~\ref{sec:tracks}. Theorem~\ref{thm:sum}
identifies the boundary contribution by two component tests; it is not
merely the residual definition~\eqref{eq:residual}. Its consequences are
proved in Section~\ref{sec:consequences}.

The closest counting framework is that of Ray and West~\cite{raywest},
whose track count is the starting point of our proof. Shar's
thesis~\cite[Corollary~4.21 and Theorem~4.25]{shar} proves eventual
polynomiality for fixed-codimension extensions of repeated direct sums
of a fixed permutation, and bounds differences between the resulting
polynomials. The examples on pp.~63--64, qualified there by an unproved
stabilization threshold, already give codimension-two predictions for
repeated $21$, $132$, $231$, and $321$. Those examples are prior work.
Our formula applies exactly to arbitrary, possibly unequal, summands,
with no stabilization condition. No equivalent composition law was found
in the Ray--West and Shar passages inspected, or in the bounded comparison
with Vatter's problem discussion, on 12 September 2026. This comparison
does not establish exhaustive priority.

\section{Insertion witnesses and good gaps}\label{sec:tracks}

We state explicitly the parts of the Ray--West framework used below.
All sites refer to the original grid of a permutation $\pi\in S_m$.
A two-point insertion triple is $t=(\rho;R,C)$, where
$\rho\in S_2$ and
\[
 R=(r_1,r_2),\quad C=(c_1,c_2),\qquad
 1\le r_1\le r_2\le m+1,\quad 1\le c_1\le c_2\le m+1.
\]
The inserted entries occupy positions $c_1,c_2+1$, and their values
are $r_1,r_2+1$, paired in relative order $\rho$. Equivalently, the
original-grid sites are $(r_{\rho(h)},c_h)$ for $h=1,2$.
The original point $(q,\pi(q))$ is sent to
\begin{equation}\label{eq:shift}
 \left(q+\#\{c\in C:c\le q\},\quad
 \pi(q)+\#\{r\in R:r\le\pi(q)\}\right),
\end{equation}
where multisets are counted with multiplicity. Write $\Ins(t;\pi)$
for the resulting permutation. These are the insertion conventions of
Ray and West~\cite[Section~4]{raywest}.

A site $(r,c)$ is \emph{active} if neither original column $c-1$ nor
original column $c$ contains value $r$; nonexistent columns are omitted.
A triple is active if its two sites are active. A \emph{synonymity}
$t\sim_\pi u$ means the equality $\Ins(t;\pi)=\Ins(u;\pi)$ of actual
permutations. Only the left triple $t$ is required to be active.
The right triple need not be active or canonical.

Put $u=(\sigma;S,D)$, with sorted multisets $S=(s_1,s_2)$ and
$D=(d_1,d_2)$. The two orientations of pattern (C) used here are
\begin{align}
 r_1\le r_2<s_1\le s_2,\quad c_1\le c_2<d_1\le d_2,
 &\qquad (+),\label{eq:plus}\\
 r_1\le r_2<s_1\le s_2,\quad d_1\le d_2<c_1\le c_2,
 &\qquad (-).\label{eq:minus}
\end{align}
The strict middle inequalities are essential; the outer inequalities
allow repeated sites. The following source facts fix the precise
meaning of the good-gap count~\cite[Proposition~6.3, Theorem~6.4,
Proposition~8.5, Corollary~8.6 and Theorem~8.7]{raywest}.

\begin{enumerate}
\item The number $j(\pi)$ counts the distinct active left triples having
a synonymous right triple in~\eqref{eq:plus} or~\eqref{eq:minus}.
There is at most one such left triple per adjacent gap. A gap is
\emph{good} if it receives a triple. Multiple right witnesses for the
same left triple are counted once.
\item The affected original points split into two nonempty tracks. In
both orientations values increase along each track; positions increase
in $(+)$ and decrease in $(-)$. The tracks cover every permutation
point of a contiguous support rectangle. Its row interval is
$[r_1,s_2-1]$, and its column interval is $[c_1,d_2-1]$ in $(+)$ or
$[d_1,c_2-1]$ in $(-)$. Each row and each column of this rectangle
has its unique permutation point inside it.
\item In $(+)$ the assigned gap is a descent, and its index $i$ satisfies
\begin{equation}\label{eq:source}
 i=c_2,\qquad r_1=\pi(i+1),\qquad \pi(c_1)=r_2+1.
\end{equation}
The points at $i,i+1$ belong to different tracks. In $(-)$ the assigned
gap is an ascent, its index is $i=c_1-2$, and its two incident points
again belong to different tracks.
\end{enumerate}

For clarity, complete support can also be read directly from the two
embeddings~\eqref{eq:shift}. The row-shift difference is positive exactly
on $[r_1,s_2-1]$. The column-shift difference has the corresponding
positive or negative sign exactly on the indicated column interval.
For an original point, one shift difference vanishes if and only if
the other does: otherwise the common ambient permutation would have
two different points in a single row or column. Thus membership in
the row interval is equivalent to membership in the column interval.
Following identifications of original points between the embeddings
gives the tracks of Theorem~6.4; Proposition~8.5 supplies their two
nonempty components. In particular, there are no omitted original
points inside the support used below.

The negative gap convention follows by horizontal reflection of the
positive case: original position $q$ becomes $m+1-q$, whereas a grid
column $c$ becomes $m+2-c$. Hence $i=c_1-2$. This explicit reflection
also removes any ambiguity caused by the repeated positive sign in
the second case of the printed proof of Corollary~8.6.
Let $\good(\pi)\subseteq\{1,\ldots,m-1\}$ denote the good-gap set.
The source count says $j(\pi)=|\good(\pi)|$.

\section{Restriction at a direct-sum cut}

Throughout this section let $\pi=\alpha\oplus\beta$ and $k=|\alpha|$.

\begin{lemma}[Restriction with re-pairing]\label{lem:repair}
Suppose $t\sim_\pi u$ has orientation $(+)$, its left triple is active,
and $\max R,\max C\le k$. Define
\[
 S'=(\min(s_1,k+1),\min(s_2,k+1)),\qquad
 D'=(\min(d_1,k+1),\min(d_2,k+1)).
\]
There exists $\sigma'\in S_2$ such that
$t\sim_\alpha(\sigma';S',D')$ is a synonymity of orientation $(+)$
with active left triple and the same assigned gap.
\end{lemma}

\begin{proof}
The insertion of $t$ has the form $\tau_A\oplus\beta$, where
$\tau_A=\Ins(t;\alpha)$ has length $k+2$. Consider the occurrence of
$\pi$ furnished by the right insertion $u$ in this common permutation.
For each of its first $k$ original points, both coordinates in
\eqref{eq:shift} are at most $k+2$. These points therefore give an
occurrence of $\alpha$ inside $\tau_A$. Its two complementary points
determine an insertion triple on $\alpha$.

For $q\le k$ and any site $d$,
\[
 \min(d,k+1)\le q\quad\Longleftrightarrow\quad d\le q.
\]
Thus $D'$ gives exactly the positions of the retained original points
in that occurrence. Those positions uniquely determine the two missing
positions, hence their original-grid insertion multiset. The same
argument on values gives $S'$. Pairing each missing position with its
actual value in $\tau_A$ determines $\sigma'$, and proves the required
insertion identity. This argument permits repeated coordinates and
clipping of either or both coordinates.

The left triple remains active upon deleting the suffix: no new
adjacent entry equal to a source row is introduced. Moreover,
$s_1>r_2$ and $k+1>r_2$ imply $\min(s_1,k+1)>r_2$.
Likewise $\min(d_1,k+1)>c_2$. Both strict inequalities survive.
The left triple and its gap index $c_2$ are unchanged.
\end{proof}

\begin{remark}\label{rem:clipping}
The right permutation $\sigma'$ cannot in general be fixed to $\sigma$.
For $\pi=213=21\oplus1$, the triples
\[
 (21;(1,1),(1,1))\quad\hbox{and}\quad(12;(4,4),(3,4))
\]
both insert to $21435$. At the cut $k=2$, coordinate clipping while
retaining $12$ gives $2134$, whereas the left insertion on $21$ is
$2143$. The repaired right triple is $(21;(3,3),(3,3))$.
This small failure is the reason restriction is formulated using the
induced occurrence and its complementary points.
\end{remark}

\begin{lemma}[Internal gaps]\label{lem:internal}
For $i<k$, $i\in\good(\pi)$ if and only if $i\in\good(\alpha)$.
For $i>k$, $i\in\good(\pi)$ if and only if $i-k\in\good(\beta)$.
\end{lemma}

\begin{proof}
First restrict a positive-orientation witness with gap $i<k$.
By~\eqref{eq:source}, $c_1\le c_2=i\le k-1$, and
$r_2=\pi(c_1)-1\le k-1$. Lemma~\ref{lem:repair} applies regardless
of the location of the right sites.
For a positive witness with $i>k$, the same equations give
$r_1=\pi(i+1)>k$ and $\pi(c_1)=r_2+1>k$, so $c_1>k$.
Strict separation puts every coordinate of $R,C,S,D$ at least $k+1$.
Subtracting $k$ from all coordinates restricts the two equal insertions
to $\beta$. Activity, separation and the translated gap are preserved.

Next consider a negative witness. A track cannot cross the direct-sum
cut: a point of $\alpha$ is both earlier and smaller than every point
of $\beta$, whereas positions and values have opposite orders along
such a track. If $i,i+1$ lie in $\alpha$, the two tracks containing
them lie wholly in $\alpha$. Since they cover the entire support,
all its rows and columns are at most $k$, and all four site multisets
lie in $[1,k+1]$. If $i,i+1$ lie in $\beta$, all support rows and
columns exceed $k$, and all sites are at least $k+1$. Restriction of
the common insertion, with translation in the latter case, proves
the required synonymity. No clipping is needed in this orientation.

Conversely, a witness on $\alpha$ lifts by leaving every site unchanged;
the common insertion becomes the old common result direct-summed with
$\beta$. A witness on $\beta$ lifts by adding $k$ to all sites.
The strict inequalities persist. For a prefix lift,
$r_2<s_1\le k+1$ forces every source row to be at most $k$.
Consequently an added adjacent entry from $\beta$ cannot spoil activity,
even at grid column $k+1$. For a suffix lift, all source rows exceed
$k$, so an adjacent prefix entry cannot spoil activity at the first
suffix grid column. This proves both converses, including boundary sites.
\end{proof}

\begin{lemma}[The join]\label{lem:join}
The gap $k$ belongs to $\good(\alpha\oplus\beta)$ if and only if
$T(\alpha)H(\beta)=1$.
\end{lemma}

\begin{proof}
The join is an ascent, so any witness there has orientation $(-)$.
Its tracks contain $\pi(k)$ and $\pi(k+1)$ separately. By the preceding
monotonicity argument, one track lies wholly in $\alpha$ and the other
wholly in $\beta$. The support rectangle has rows $[r_1,s_2-1]$ and
columns $[d_1,c_2-1]$. Its prefix part occupies exactly columns
$d_1,\ldots,k$ and values $r_1,\ldots,k$. Every such row and column
is occupied, so counting gives $d_1=r_1$. All these points lie on one
decreasing track; hence
\[
 (\pi(d_1),\ldots,\pi(k))=(k,k-1,\ldots,d_1).
\]
Similarly, its suffix part is a decreasing initial interval of $\beta$,
using values $k+1,\ldots,c_2-1$; counting gives $s_2=c_2$.
The complementary prefix entries use smaller values, and the remaining
suffix entries use larger values. These intervals are therefore
direct-sum components. A decreasing permutation is sum-indecomposable,
including in size one, so they are precisely the two incident maximal
components. This proves necessity for arbitrary summands.

For sufficiency, write $\delta_a=(a,a-1,\ldots,1)$.
It is enough to exhibit a witness on $\delta_a\oplus\delta_b$ for
all $a,b\ge1$, then lift across the surrounding direct summands using
Lemma~\ref{lem:internal}. Put $m=a+b$ and take
\begin{equation}\label{eq:witness}
 \rho=\sigma=12,\quad R=D=(1,a),\quad C=S=(a+2,m+1).
\end{equation}
By~\eqref{eq:shift}, both insertions equal
\begin{equation}\label{eq:ambient}
 (a+2,\ a,a-1,\ldots,2,\ m+2,\ 1,\ m+1,m,\ldots,a+3,\ a+1),
\end{equation}
with empty decreasing ranges omitted. At source site $(1,a+2)$
the adjacent entries are greater than $1$; when $b=1$, only the
adjacent entry $a+1$ exists. At $(a,m+1)$ the only adjacent entry is
$a+1$. The left triple is active, and both strict inequalities reduce
to $a<a+2$. Its gap is $c_1-2=a$. The construction includes $a=1$
or $b=1$; when $a=b=1$ the common ambient permutation is $3412$.
\end{proof}

\begin{proof}[Proof of Theorem~\ref{thm:sum}]
Lemma~\ref{lem:internal} identifies all internal good gaps, and
Lemma~\ref{lem:join} identifies the only remaining one. In fact,
\[
 \good(\alpha\oplus\beta)=\good(\alpha)\ \cup\
 \{k+i:i\in\good(\beta)\}\ \cup\
 \begin{cases}\{k\},&T(\alpha)H(\beta)=1,\\
 \varnothing,&T(\alpha)H(\beta)=0.
 \end{cases}
\]
The sets are disjoint. Counting them with $j=|\good|$ proves
\eqref{eq:sum} and the stronger gap statement.
\end{proof}

\section{Skew sums, separable permutations and equality}\label{sec:consequences}

Let $c(\pi)(q)=|\pi|+1-\pi(q)$ be value complementation. It gives a
bijection between ambient permutations containing $\pi$ and those
containing $c(\pi)$, so
\begin{equation}\label{eq:complement}
 j(c(\pi))=j(\pi),\qquad
 c(\alpha\skewsum\beta)=c(\alpha)\oplus c(\beta).
\end{equation}
Let $H_-(\pi),T_-(\pi)$ indicate whether the first and last
skew-indecomposable components, respectively, are increasing.

\begin{corollary}\label{cor:skew}
For all nonempty $\alpha,\beta$,
\[
 j(\alpha\skewsum\beta)=j(\alpha)+j(\beta)+T_-(\alpha)H_-(\beta).
\]
\end{corollary}
\begin{proof}
Complementation converts maximal skew components to maximal sum
components and increasing permutations to decreasing permutations.
Apply Theorem~\ref{thm:sum} to~\eqref{eq:complement}.
\end{proof}

For a maximal sum decomposition $\pi=\pi_1\oplus\cdots\oplus\pi_t$,
iteration of Theorem~\ref{thm:sum} yields
\begin{equation}\label{eq:maximal}
 j(\pi)=\sum_{h=1}^{t}j(\pi_h)
 +\sum_{h=1}^{t-1}\ind\{\pi_h,\pi_{h+1}\text{ are decreasing}\}.
\end{equation}
The maximal skew version replaces ``decreasing'' by ``increasing''.

A \emph{separable} permutation is obtained from singletons by repeated
direct and skew sums, as in Albert, Atkinson and
Vatter~\cite[Section~1]{albert}. Recursively taking maximal decompositions
gives its ordered decomposition tree: each internal node is labelled
$\oplus$ or $\skewsum$ and its children are its maximal components.
For length at least two the two decomposability types cannot occur
simultaneously. To see this, a proper initial segment cannot be both
the smallest and the largest set of values, and two proposed initial
segments of these different types are nested. Thus the type at each
node is unambiguous. Every leaf is a singleton.

\begin{corollary}[An intrinsic statistic]\label{cor:tree}
At each sum node count adjacent children whose represented permutations
are both decreasing, and at each skew node count adjacent children whose
represented permutations are both increasing. The total over all nodes
of the maximal decomposition tree of a separable permutation $\pi$ is
$j(\pi)$.
\end{corollary}
\begin{proof}
For the singleton, $g_2(1)=6=B(1)$ and $j(1)=0$.
Apply~\eqref{eq:maximal} or its skew version at each internal node,
and induct on the tree. The terms at singleton leaves vanish.
\end{proof}

This statistic can be evaluated from the permutation's decomposition,
without enumerating insertion triples or ambient containing permutations.
It also determines $g_2$ by~\eqref{eq:residual}.

A permutation is \emph{layered} if it is a direct sum of decreasing
permutations, and \emph{colayered} if it is a skew sum of increasing
permutations.

\begin{corollary}[Equality on decomposable permutations]\label{cor:equality}
If $\pi$ is sum-decomposable, then
$j(\pi)=|\pi|-1$ if and only if $\pi$ is layered.
If $\pi$ is skew-decomposable, then
$j(\pi)=|\pi|-1$ if and only if $\pi$ is colayered.
In particular, among separable permutations equality holds precisely
for the layered and colayered permutations.
\end{corollary}
\begin{proof}
For the maximal sum components $\pi_1,\ldots,\pi_t$, with $t\ge2$,
let $I_h$ indicate that $\pi_h$ and $\pi_{h+1}$ are both decreasing.
Equation~\eqref{eq:maximal} gives the deficit identity
\begin{equation}\label{eq:deficit}
 |\pi|-1-j(\pi)=
 \sum_{h=1}^{t}\bigl(|\pi_h|-1-j(\pi_h)\bigr)
 +\sum_{h=1}^{t-1}(1-I_h).
\end{equation}
Each term is nonnegative by the Ray--West bound. Equality forces every
$I_h=1$. Since $t\ge2$, each child is in an adjacent pair and therefore
is decreasing. Hence $\pi$ is layered. Conversely, the skew law on
singletons gives $j(\delta_a)=a-1$, so decreasing children and
\eqref{eq:maximal} imply equality for every layered sum-decomposable
permutation. Complementation proves the skew statement. Every separable
permutation of size at least two has a decomposition of one type; in
size one equality and both descriptions hold directly.
\end{proof}

\section{Exact checks and remaining questions}

The proof above is uniform in the lengths of the summands. Exact finite
computations served to test its conventions and boundary cases.
An insertion-based verifier examined every permutation through length
$6$ and $15$ structured examples of lengths $10$--$12$. It generated
all two-point insertion triples, permitting inactive right triples,
reconstructed the tracks from the two embeddings, and checked support,
source equations, internal restriction and the join criterion.
It tested $3{,}110$ witness/cut cases; $1{,}219$ used positive-prefix
restriction and $188$ required changing the right inserted permutation.
All checks passed.

A separately authored verifier reconstructed the calculation from
ambient permutations and two-position deletions, without importing
the insertion generator. This separation of generation from checking
is also used concretely by Zeng, Liu, Ratnavelu and
Ong~\cite[Section~6]{zengcert}, whose finite-geometry certificates are
checked afresh by a standalone verifier. Here the second program
enumerated all $46{,}230$ ambient permutations of lengths $3$--$8$,
and all $1{,}246{,}962$ two-position deletion occurrences. Deduplication
by ambient permutation gave $g_2$ for all $873$ patterns of lengths
$1$--$6$. Besides agreement with the first counts, it verified all
$465$ direct-sum cuts and their exact good-gap identities. Its $659$
positive-prefix restrictions included $134$ changes of right pattern.
A separate deletion-decoding check verified~\eqref{eq:witness}--\eqref{eq:ambient}
for every $1\le a,b\le50$, including repeated and boundary sites.

The tree implementation agreed with the inherited exact census for all
$52{,}465$ separable permutations through length $9$. Table~\ref{tab:coverage}
records that range. The second verifier separately confirmed the $515$
separable permutations through length $6$. These are exact finite
checks and implementation comparisons, not proofs for untested lengths.
The inherited census informed the statement; it is exploratory and
regression evidence, not a blind experiment. Separate code and a fresh
review context do not make previously seen data statistically independent.

\begin{table}[ht]
\centering
\caption{Coverage of the inherited census used for the tree regression.
Every listed separable permutation agrees with Corollary~\ref{cor:tree}.}
\label{tab:coverage}
\small
\begin{tabular}{lrrrrrrrrr}
\toprule
Length & 1 & 2 & 3 & 4 & 5 & 6 & 7 & 8 & 9\\
\midrule
All patterns & 1 & 2 & 6 & 24 & 120 & 720 & 5040 & 40320 & 362880\\
Separable & 1 & 2 & 6 & 22 & 90 & 394 & 1806 & 8558 & 41586\\
\bottomrule
\end{tabular}
\end{table}

The remaining issue in Vatter's problem is an intrinsic evaluation of
$j$ when repeated sum and skew decomposition leaves a permutation
indecomposable in both senses. The equality classification proved here
also leaves that domain untreated. Neither a general deletion inequality
for $j$ nor a general equality classification is needed in our argument.
The natural next test of any proposed extension is an exact counterexample
search on permutations indecomposable in both senses, followed by a
uniform proof if the proposed rule survives.

\paragraph{Computation and assistance disclosure.}
The companion reproducibility instructions identify the programs,
coverage reports and exact census inputs. The counting and witness
checks use integer arithmetic; the census array is stored as exact
integers. No numerical fitting or tolerance enters the stated identities.
OpenAI Codex assisted with literature searches, proof development and
revision, program construction, a separate mathematical review pass,
and manuscript preparation. That review is a separate model context,
not external peer review. Language-model output and successful tests
are not substitutes for the written proof. No proof assistant was used.

\bibliographystyle{amsplain}
\bibliography{references}
\end{document}
